\documentclass{article}

% NeurIPS 2025 style
\usepackage[final,nonatbib]{neurips_2025}

\usepackage[utf8]{inputenc}
\usepackage[T1]{fontenc}
\usepackage{hyperref}
\usepackage{url}
\usepackage{booktabs}
\usepackage{amsfonts}
\usepackage{amsmath}
\usepackage{nicefrac}
\usepackage{microtype}
\usepackage{xcolor}

\title{Cold-Start Predictive Evaluation via\\Rasch-NCF Pipeline with Online Platt Calibration}

\author{
  Lisa St John \\
  CS321M: AI Measurement Science \\
  \texttt{lisa.stephan.sf@gmail.com}
}

\begin{document}

\maketitle

\begin{abstract}
We describe our approach to the CS321M Predictive Evaluation Challenge.
We adopt a three-stage pipeline:
(1)~a Rasch IRT model estimates subject ability $\theta$ and item difficulty $\beta$ via Joint MLE;
(2)~a content-based Neural Collaborative Filtering (NCF) model maps text embeddings from a Sentence Transformer (\texttt{all-mpnet-base-v2}, 768-dim) together with a learned condition embedding (216 conditions $\times$ 32-dim) through an MLP head to produce logits for cold-start items;
(3)~online Platt temperature scaling calibrates probabilities using the small set of adaptively labeled examples each round.
On 3-fold cold-start cross-validation, our method achieves Accuracy $0.800 \pm 0.034$, AUC-ROC $0.900 \pm 0.015$, with NLL coefficient of variation 3.2\%.
We also implement a farthest-point diversity acquisition function and analyze failure modes across benchmarks.
\end{abstract}

\section{Problem Formulation and Method Description}

\subsection{Problem Definition}

The Predictive Evaluation Challenge asks: given an AI system (subject) and a test question (item), predict the probability that the subject answers correctly---\emph{without} running the evaluation.
Formally, for subject $s$, item $i$, benchmark $b$, and test condition $c$:
\begin{equation}
  \hat{p} = P(Y_{si} = 1 \mid \text{subject\_content}, \text{item\_content}, b, c)
\end{equation}
All test items are cold-start (unseen during training), while subjects appear in the training set.
The primary metric is negative log-loss (higher is better).

\subsection{Three-Stage Pipeline}

\paragraph{Stage 1: Rasch IRT Parameter Estimation.}
We use the Rasch model $P(Y=1|\theta,\beta) = \sigma(\theta - \beta)$ where $\theta$ is subject ability, $\beta$ is item difficulty, and $\sigma$ is the sigmoid function.
Parameters are estimated via Joint Maximum Likelihood Estimation (JMLE) using Adam (lr$=$0.01) for 50 epochs with $L_2$ regularization ($\lambda=10^{-4}$) and ReduceLROnPlateau scheduling.
Training accuracy reaches 0.838.

\paragraph{Stage 2: Content-Based NCF.}
Since test items are cold-start, we learn a mapping from text content to prediction logits.
A Sentence Transformer (\texttt{all-mpnet-base-v2}) encodes \texttt{subject\_content} and \texttt{item\_content} into $L_2$-normalized 768-dim vectors.
A condition embedding table (216 conditions $\times$ 32-dim, with an ``unknown'' fallback for unseen conditions) captures test-condition effects.
The MLP head takes the concatenation $[u; v; c_\text{emb}] \in \mathbb{R}^{1568}$ through two hidden layers (256 units, BatchNorm, ReLU, Dropout) to produce a scalar logit.

Training uses BCEWithLogitsLoss, Adam (lr$=5 \times 10^{-4}$, weight decay $10^{-4}$), gradient clipping (max norm 1.0), batch size 4096, up to 50 epochs with early stopping (patience$=$8).
Model selection uses \textbf{validation NLL} (not accuracy) to align with the competition metric.

At inference, an LRU embedding cache (capacity 2000) avoids redundant Sentence Transformer calls---since the same subject repeats across all items in a round, caching reduces encoding calls by ${\sim}50\%$ and prevents Codabench timeouts.

\paragraph{Stage 3: Online Platt Calibration.}
Each round, the platform reveals $K{=}5$ labeled examples per data category via adaptive labeling.
We fit a temperature parameter $T^*$ by grid search over $T \in [0.2, 5.0]$ (200 points), minimizing NLL on the labeled set:
$p_\text{cal} = \sigma(\text{logit} / T^*)$.
A safety guard ensures we only use $T^* \neq 1$ when calibration strictly improves NLL; otherwise we fall back to $T{=}1.0$.
The calibrator auto-resets when the labeled set changes between rounds.
Final probabilities are clipped to $[0.01, 0.99]$.
On any exception, \texttt{predict()} falls back to the labeled base rate.

\paragraph{Acquisition Function.}
We implement greedy farthest-point diversity sampling: the first candidate seeds the pool, and subsequent candidates score by their minimum $L_2$ distance to all previously seen embeddings.
This encourages label diversity without requiring pre-fit centroids.
A \texttt{math.isfinite} guard prevents non-finite returns from triggering the platform's random-selection fallback.

\section{Training Data Description}

We use the public \texttt{aims-foundations/measurement-db} dataset from HuggingFace.
The raw dataset contains 10.7M response rows across 909 subjects and 103,983 items spanning multiple benchmarks.

\paragraph{Key preprocessing.}
Approximately 8.5\% of responses are non-binary floats (from rating-scale benchmarks such as \texttt{ultrafeedback} and \texttt{mt\_bench}).
Since both the Rasch model and BCEWithLogitsLoss require binary labels, we \emph{first} filter to $\text{response} \in \{0, 1\}$, \emph{then} cast to int---this ordering prevents silent truncation (e.g., 0.7 $\to$ 0).
The filtered dataset contains \textbf{4,448,512 rows}, \textbf{70,873 items}, and \textbf{216 test conditions}, with an overall pass rate of ${\sim}0.63$ and response-matrix sparsity of ${\sim}93\%$.

The data exhibits a long-tail distribution: a few large benchmarks (MMLU, GSM8K) contribute most rows, while many small benchmarks have few items.
Item difficulty is bimodal---many items are very easy ($>$0.95 pass rate) or very hard ($<$0.05), with relatively few at moderate difficulty.

\section{Experimental Results and Ablation Studies}

\subsection{Cold-Start Cross-Validation}

We evaluate with item-stratified $K$-fold CV simulating the cold-start regime: test items never appear in training, while subjects are shared.
Each fold uses 25 adaptive labels for Platt calibration.

\begin{table}[h]
\centering
\caption{3-fold cold-start cross-validation results.}
\label{tab:cv}
\begin{tabular}{lcccc}
\toprule
Fold & Neg-NLL & Accuracy & AUC-ROC & vs.\ Baseline \\
\midrule
1 & $-0.414$ & 0.830 & 0.919 & $-33.9\%$ \\
2 & $-0.393$ & 0.820 & 0.897 & $-36.6\%$ \\
3 & $-0.424$ & 0.750 & 0.884 & $-36.2\%$ \\
\midrule
Mean$\pm$Std & $-0.410 \pm 0.013$ & $0.800 \pm 0.034$ & $0.900 \pm 0.015$ & $-35.6\%$ \\
\bottomrule
\end{tabular}
\end{table}

The model consistently outperforms the base-rate baseline (predicting the global pass rate) across all folds, with NLL improvement averaging 35.6\%.
The NLL coefficient of variation is 3.2\% ($<$5\%), indicating stable predictions across different item subsets.

\subsection{Per-Benchmark Performance}

\begin{table}[h]
\centering
\caption{Per-benchmark performance (mean$\pm$std across 3 folds).}
\label{tab:bench}
\begin{tabular}{lccc}
\toprule
Benchmark & Neg-NLL & Accuracy & AUC-ROC \\
\midrule
livecodebench & $-0.179 \pm 0.084$ & $0.958 \pm 0.059$ & $1.000 \pm 0.000$ \\
ai2d\_test & $-0.450 \pm 0.120$ & $0.791 \pm 0.048$ & $0.865 \pm 0.040$ \\
bfcl & $-0.496 \pm 0.127$ & $0.712 \pm 0.052$ & $0.866 \pm 0.066$ \\
mathvista\_mini & $-0.551 \pm 0.149$ & $0.750 \pm 0.083$ & $0.738 \pm 0.138$ \\
mmbench\_v11 & $-0.293 \pm 0.035$ & $0.868 \pm 0.079$ & $0.684 \pm 0.224$ \\
mlupro & $-0.553 \pm 0.188$ & $0.729 \pm 0.045$ & $0.886 \pm 0.064$ \\
rewardbench & $-0.514 \pm 0.098$ & $0.741 \pm 0.104$ & $0.766 \pm 0.104$ \\
\bottomrule
\end{tabular}
\end{table}

\texttt{livecodebench} achieves the best performance (AUC 1.000), likely because coding tasks have clear difficulty gradients that the encoder captures well.
\texttt{rewardbench} performs worst (AUC 0.766), likely because reward-model evaluation involves subtle semantic distinctions that a general-purpose encoder struggles with.

\subsection{Ablation Studies}

\paragraph{Rasch baseline vs.\ NCF.}
The Rasch model (Stage~1) achieves training accuracy 0.838 but cannot handle cold-start items.
NCF (Stage~2) lifts cold-start accuracy to 0.847 on the training validation split and 0.800 in cross-validation, confirming effective text-to-parameter generalization.

\paragraph{Condition embedding.}
The 216-condition embedding table lets the model distinguish the same item under different test conditions (e.g., zero-shot vs.\ chain-of-thought), capturing condition-specific performance variation.

\paragraph{Platt calibration.}
Online calibration with 25 labeled examples fits $T^*$ via grid search.
The safety fallback ($T{=}1$ when calibration hurts) ensures calibration never degrades performance, though improvement is modest with few labels.

\paragraph{Model selection by NLL.}
Selecting checkpoints by validation NLL (instead of accuracy) directly aligns with the competition metric, avoiding the risk of choosing high-accuracy but poorly-calibrated models.

\paragraph{Embedding cache.}
The LRU cache (capacity 2000) reduces total encoding calls by ${\sim}50\%$ per round, critical for meeting the per-round time budget on Codabench.

\paragraph{Adaptive vs.\ random labeling.}
Farthest-point diversity acquisition shows only a 0.0006 NLL difference from random selection.
Under the current single-temperature Platt framework, \emph{which} examples are labeled matters less than having \emph{some} labeled examples.
Per-benchmark calibration or uncertainty-aware acquisition could improve this.

\section{Failure Mode Analysis}

\paragraph{Model stability.}
Early 5-fold CV showed NLL CV of 14.4\%, indicating unstable predictions.
After widening the temperature search range, adding auto-reset, and enabling embedding caching, NLL CV dropped to 3.2\%.
Residual variance comes from benchmark-level imbalance: large benchmarks provide strong training signal while small ones generalize poorly.

\paragraph{Non-binary responses.}
About 8.5\% of raw responses are non-binary floats from rating-scale benchmarks (\texttt{ultrafeedback}, \texttt{mt\_bench}).
These caused loss explosion ($-2.5 \times 10^9 \to -7.5 \times 10^{12}$) before filtering.
The filter-then-cast ordering is critical; reversed ordering silently truncates 0.7 to 0.
These benchmarks are fully excluded post-filtering---if the test set includes their items, the model lacks targeted training signal.

\paragraph{Encoder limitations.}
\texttt{all-mpnet-base-v2} is a general-purpose text encoder.
For highly structured content (mathematical formulas, code snippets), encoding quality may degrade.
The lower AUC on \texttt{rewardbench} (0.766) likely reflects this: reward-model evaluation requires fine-grained semantic discrimination beyond the encoder's capacity.

\paragraph{Fundamental cold-start limits.}
Our method infers item difficulty from text alone, but difficulty often depends on subtle details invisible in the text.
Two similar-looking math problems may differ drastically in difficulty due to a small numerical change---a fundamental limitation of pure content-based approaches.

\paragraph{Future directions.}
(1)~2PL/3PL IRT models with discrimination and guessing parameters;
(2)~cross-benchmark transfer to improve small-benchmark predictions;
(3)~uncertainty-aware acquisition combining predictive entropy with diversity;
(4)~domain-specific encoders for code and mathematics;
(5)~per-benchmark temperature calibration instead of a single global $T$;
(6)~LayerNorm to eliminate BatchNorm's reliance on running statistics at inference.

\end{document}
